\section{Research question}
\label{sec:question}

The literature does not test provenance verification, independent rebuilding, and DNS monitoring together in a SolarWinds-style lab, and the PoC implements none of them as defences. That gap gives the question:

\begin{quote}
\textbf{How effectively can independent build-artifact verification and DNS activity monitoring detect different stages of a compromised but legitimately signed software-update workflow in a controlled Mininet/Docker simulation?}
\end{quote}

Three sub-questions separate the stages. Before release, can an independently generated reference artifact or a trusted provenance policy distinguish the tampered V3 from an approved release when both pass signature checks? After installation, can DNS logging reveal the check-in and CNAME handoff, and can a stated rule separate them from benign lab traffic? And which stages remain invisible when either control is used alone?

The evaluation would run a baseline approved update, the tampered update, an artifact comparison, DNS monitoring on labelled runs, and both controls together, recording detected releases, flagged check-ins, false positives on benign runs, and the stage at which each alert fires. The lab would need benign comparison traffic before any detection rate means anything. If reproducible builds cannot be implemented, an approved-reference comparison is the narrower fallback, and it must not be described as proof of reproducible-build security.

The limits are the lab's: a predefined check-in routine, simplified DNS, no SUNSPOT replacement, no original DGA, and no later-stage malware. Results would describe this environment, not SolarWinds.
